%% ma: L11-B10-0006
%% lop: 11
%% chuong: 4
%% bai: 10
%% dang: 11.10.3
%% loai: TN4
%% muc: TH
%% dap_an: "A"
%% nguon: "de-cuong-giua-ky-1-toan-11-nam-2026-2027-truong-thpt-hoang-van-thu-ha-noi-2-8.pdf (M08, Đề cương giữa kì 1 Toán 11 THPT Hoàng Văn Thụ, Hà Nội), Đề minh họa, Phần I Câu 2"
%% kien_thuc: "Bài 10 (giao điểm của đường thẳng và mặt phẳng)"
%% trang_thai: chinh-thuc
%% ghi_chu: "Đã duyệt 29/10/2026."
\begin{ex}
Cho bốn điểm $A$, $B$, $C$, $D$ không đồng phẳng. Gọi $M$, $N$ lần lượt là trung điểm của $AC$ và $BC$. Trên đoạn $BD$ lấy điểm $P$ sao cho $BP=2PD$. Giao điểm của đường thẳng $NP$ và mặt phẳng $(ACD)$ là giao điểm của
\choice{\True $NP$ và $CD$}{$NP$ và $AB$}{$NP$ và $AC$}{$NP$ và $AD$}
\loigiai{$NP\subset(BCD)$ và $(BCD)\cap(ACD)=CD$. Vì $\dfrac{BN}{BC}=\dfrac12\ne\dfrac{BP}{BD}=\dfrac23$ nên $NP$ cắt $CD$; giao điểm đó là giao điểm của $NP$ với $(ACD)$.}
\end{ex}
